\documentclass[11pt]{article}
\usepackage[a4paper,margin=26mm]{geometry}
\usepackage{amsmath,amssymb,amsthm,mathtools}
\usepackage[T1]{fontenc}
\usepackage{lmodern,microtype}
\usepackage[colorlinks=true,linkcolor=blue,urlcolor=blue,citecolor=blue]{hyperref}
\usepackage{enumitem,booktabs,array}
\setlist{itemsep=3pt,topsep=5pt}
\newtheorem{theorem}{Theorem}
\newtheorem{proposition}[theorem]{Proposition}
\newtheorem{lemma}[theorem]{Lemma}
\newcommand{\E}{\mathbb E}
\newcommand{\Prb}{\mathbb P}
\newcommand{\F}{\mathbb F}
\newcommand{\eps}{\varepsilon}
\newcommand{\rank}{\operatorname{rank}}
\newcommand{\Hom}{\operatorname{Hom}}
\newcommand{\Tr}{\operatorname{Tr}}
\newcommand{\Span}{\operatorname{span}}
\newcommand{\val}{\operatorname{val}}
\newcommand{\OPT}{\operatorname{OPT}}
\newcommand{\ind}{\mathbf 1}
\title{A Compressed Presentation of OpenAI's Proof\\ of the Unique Games Theorem}
\author{Ariel Elboim\\\small Independent researcher, ORCID 0009-0001-1341-366X\\\small Prepared with Codex (OpenAI) and Claude (Anthropic)}
\date{9 October 2026}
\newcommand{\depositdoi}{10.5281/zenodo.23264902}
\begin{document}
\maketitle
\begin{abstract}
We reconstruct the main argument of OpenAI's preprint \emph{The Unique
Games Theorem} \cite{UGT}. Three earlier theorems are stated as inputs.
The remaining argument constructs stable nonlinear noise, extracts
constant-output certificates, decodes them using only local information,
and contradicts parallel repetition on clean coordinates. We retain the
conditioning and parameter order on which these steps depend. The aim is
a short, continuous proof for a reader comfortable with finite-field
linear algebra and Fourier analysis on finite groups who has not read the
preprint. We modify the root construction of the noise gadget, simplify
certificate extraction, and rewrite the genericity proof; these changes
are identified and proved explicitly, and are listed, with smaller ones,
under Attribution.
\end{abstract}

\noindent\textbf{Status.}
This is an explanatory reconstruction. Its mathematical content
received source-based reviews by AI agents (separate Claude instances and
Codex); Codex had also helped write an earlier version. These reviews do
not certify the preprint or constitute formal verification of this note,
and no human expert has reviewed it. OpenAI's repository also contains a
Lean~4 formalization of the source theorem, whose Comparator check
passed when we ran it; that check concerns OpenAI's formal development,
not this note or its changes (Appendix~\ref{app:lean}). This note is not
derived from the formalization. Fortnow's earlier exposition
\cite{Fortnow} presents the same proof at greater length; this note was
prepared independently of it. A shorter proof, which simplifies several
of the steps presented here, is given in \cite{ShortProof}. The theorem below is the preprint's claim,
whose proof argument we present. The root step, Lemma~\ref{lem:root},
has no counterpart in the source; it and our other departures from the
source are listed under Attribution, at the end of
Section~\ref{sec:finish}.

\section{The claim and three prior inputs}

A unique game labels vertices from a finite alphabet. Each edge carries
a permutation: a label at one endpoint specifies exactly one acceptable
label at the other. Its value is the largest fraction of constraints
satisfied by one labeling.

\begin{theorem}[Main claim of \cite{UGT}]\label{thm:main}
For every fixed $0<\eps,\delta<1/2$, there is a deterministic polynomial-time
reduction from satisfiability to explicit unweighted unique games on
simple bipartite graphs that maps satisfiable formulas to YES games and
unsatisfiable formulas to NO games, where
\[
 \mathrm{YES}:\val(G)\ge1-\eps,
 \qquad \mathrm{NO}:\val(G)\le\delta.
\]
The alphabet is a fixed binary vector space, and every constraint is a
translation. The alphabet and the degree of the running-time polynomial may depend
on the two target errors, but not on the input length.
\end{theorem}

All spaces and linear maps below are over $\F_2$, except where a larger
finite field is specified. Uniform samples include zero. Separately
specified random choices are independent unless stated otherwise.
Intermediate games may have rational weights and repeated constraint
occurrences.

\paragraph{Input 1: the parity gap.}
For every fixed positive rational $\xi$, H\aa stad's theorem
\cite[Theorem 5.4 and its proof]{Hastad} gives a deterministic
polynomial-time reduction from satisfiability to weighted three-bit equations
\[
 x_{v_1}+x_{v_2}+x_{v_3}=b_e,
 \qquad \mathrm{YES}:\OPT\ge1-\xi,
 \qquad \mathrm{NO}:\OPT\le\tfrac12+\xi.
\]
The rational weights have polynomial encoding length.

\paragraph{Input 2: Grassmann expansion.}
The binary Grassmann graph has the $\ell$-dimensional subspaces of an
ambient space as vertices; adjacency means an intersection of dimension
$\ell-1$. The retention of a nonempty set is the probability that a uniform
neighbor of a uniform member stays in the set. For every fixed retention
threshold $\zeta>0$, there are $\alpha>0$ and an integer $r_{\rm s}\ge1$, depending
only on $\zeta$, such that every set with retention at least $\zeta$ has
relative density at least $\alpha$ in an interval
\[
 \{L:A_0\subseteq L\subseteq B_0\},
 \qquad \dim A_0+\operatorname{codim}B_0\le r_{\rm s}.
\]
This holds for every sufficiently large $\ell$, and then every sufficiently
large ambient dimension. This is the expansion theorem of
Khot--Minzer--Safra \cite[Theorem 1.12]{KMS}.

\paragraph{Input 3: projection-game repetition.}
In a projection game, an answer of the first player determines one required
answer of the second. Dinur--Steurer \cite[Corollary 1.2]{DS} give
\begin{equation}\label{eq:repetition}
 \val(G)\le1-g\quad\Longrightarrow\quad
 \val(G^{\otimes n})\le(1-g^2/16)^n.
\end{equation}
This applies to finite weighted games and independent constraint samples.
Each player may use its entire question tuple. The bound does not depend
on the alphabet. Our games have a fixed projection for each question pair;
repeated samples of the same constraint simply contribute their weights.

\paragraph{Idea and plan.}
The honest labeling applies one nonlinear map $C$, and the test perturbs
its input by a special noise. The noise rarely changes the output of $C$,
yet every linear family that sees enough of the output notices it with
constant probability. Section~\ref{latent:section} builds such noise by a
race inside a recursive construction. At each level an output change
survives with probability $1-\theta$, so after $s$ levels it survives with
probability about $e^{-\theta s}$. A generic family of linear observations
of rank $r$ loses at most one dimension per level, and does so with
probability at most $3r\theta$. Like a coupon collector, it therefore needs
at least about $H_r/(3\theta)$ levels to die, where
$H_r=1+\frac12+\dots+\frac1r$. The harmonic numbers diverge, so a large
initial rank wins; the proof takes $s\approx H_{r_0}/(8\theta)$. The algebra
needed is that, for a generic family, at most $3r$ of the $\theta^{-1}$ field
lines are bad, each costing at most one dimension. A bad-line bound such as
$2^r$ would make the corresponding sum converge, and this argument would
fail. Stability gives completeness
(Section~\ref{sec:tables}). Conversely, detection and Input~2 turn a labeling
that passes the test into constant-output certificates, and these give two
local decoders with constant agreement (Section~\ref{sec:decoders}). In a NO
instance, clean coordinates embed the repeated base game, and
\eqref{eq:repetition} makes such agreement impossible
(Section~\ref{sec:clean}). Section~\ref{sec:finish} removes weights and
amplifies the gap.

\paragraph{Proof map.}
Four statements carry the argument; the notation is defined in
Sections~2--5.
\begin{enumerate}[label=(\roman*)]
\item \emph{Gadget:} the contract \eqref{eq:gadget-contract}, with $r_*$
 depending only on $p_*$ (Section~\ref{latent:section}, Lemma~\ref{lem:root}).
\item \emph{Completeness:} \eqref{eq:completeness}.
\item \emph{Extraction:} for all cubes $k$ beyond a threshold, any labeling
 that passes the test with probability at least $99/100$ yields two local
 strategies with legal answers that agree with probability at least a
 constant $\gamma>0$ (Proposition~\ref{prop:extraction}).
\item \emph{Clean coordinates:} if $\OPT(\mathcal E)\le4/5$, as in a NO
 instance with $\xi<1/100$, every pair of local
 strategies with legal answers in the same experiment agrees with
 probability at most $\exp[-2^{-\dim W-r_{\rm s}}k^{1/3}/3600]$
 (Proposition~\ref{prop:clean}).
\end{enumerate}
The threshold, $\gamma$, $\dim W$ and $r_{\rm s}$ are fixed before $k$ and do not
depend on the instance or the labeling; this matters because the parity
error, and hence the instance, is chosen after $k$. So (iii) and (iv) give
the NO side of the gap \eqref{eq:intermediate-gap}, and (ii) gives its YES
side. Section~\ref{sec:finish} turns the gap
into Theorem~\ref{thm:main}.

\section{From parity equations to shared tables}\label{sec:tables}

\paragraph{Distinct variables in each equation.}
Let $\mathcal E_0$ be the system given by Input~1. Replace every variable by
$100$ clones. In each equation, enumerate uniform
clone choices conditioned on repeated occurrences of one variable using
distinct clones. Split the original weight among these choices. Copying
a global assignment to all clones preserves completeness.

For an arbitrary clone assignment put $t_v=\E_c(-1)^{x_{v,c}}$. Before
conditioning, the equation's satisfaction probability is
\[
 \frac{1+(-1)^{b_e}t_{v_1}t_{v_2}t_{v_3}}2.
\]
Round each variable by the sign of its bias, breaking zero ties arbitrarily.
An equation violated after rounding has satisfaction probability at most
$1/2$ before conditioning; any other equation has probability at most one.
The probability of a forbidden clone collision is at most $3/100$.
Conditioning away an event of probability $p$ changes the expectation of
a $[0,1]$-valued variable by at most $p$. Write $\mathcal E$ for the
processed system and $\omega_e$ for its weights. Then
$\OPT(\mathcal E)\ge\OPT(\mathcal E_0)$ and
\[
 \OPT(\mathcal E)\le\frac{1+\OPT(\mathcal E_0)}2+\frac3{100},
\]
so in the NO case, for $\xi<1/100$, $\OPT(\mathcal E)\le4/5$.

In its equation--variable game, the first player receives an equation and
returns a satisfying triple. The second receives a uniformly chosen one
of its three variables and returns a bit. The second player's answers form
a global assignment. On an equation this assignment violates, every legal
triple disagrees with it in at least one position. Hence in the NO case
\begin{equation}\label{eq:base-gap}
 \val(G_{\mathcal E})\le1-\frac{1-\OPT(\mathcal E)}3\le14/15.
\end{equation}
Distinct variables ensure that an equation--variable question pair specifies
a unique position.

\paragraph{The noise gadget needed below.}
Section~\ref{latent:section} constructs, for each $p_*>0$, a threshold $r_*$
such that for every sufficiently large $\ell$ divisible by $3$ there are
finite data
\begin{equation}\label{eq:gadget-contract}
 \begin{gathered}
 K=\F_2^\ell\le\mathcal V,\quad C:\mathcal V\to K,\quad a\sim\mu,\\
 C(x+h)=C(x)+h\quad(h\in K),\qquad
 \Pr_{x\text{ uniform},a}[C(x+a)\ne C(x)]\le p_*,\\
 G\le\mathcal V^*,\ \dim(G|_K)\ge r_*
 \quad\Longrightarrow\quad
 \Pr_a[\exists g\in G:g(a)\ne0]\ge1/8.
 \end{gathered}
\end{equation}
The noise law is rational. Crucially, $r_*$ is chosen before $\ell$.
Fix a linear splitting $\mathcal V=K\oplus W$.

\paragraph{Questions and homogeneous answers.}
Sample $U=(e_1,\ldots,e_k)$ independently with replacement. Put
$L_U=\F_2^{1+2k}$, with coordinates $t_0,x_1,y_1,\ldots,x_k,y_k$.
The triple for equation $e_i$ is
\[
 (x_i,y_i,b_{e_i}t_0+x_i+y_i).
\]
Vectors with $t_0=1$ are exactly the legal answers. Different positions
retain separate answer coordinates, even when their variable identifiers
repeat. A projected question $O$ retains either an equation or one of its
actual bits in each position; a singleton position records only the
variable's identifier, not the equation or the slot it came from. This gives a surjective linear map
$\pi:L_U\to L_O$ preserving $t_0$.

\paragraph{Table keys force sharing.}
A table is a linear map $P:L_Q\to\mathcal V$, or an affine map on the
legal answers. In each position a table depends on the whole equation, on
one of its variables, or on nothing. Its key is the coarsest support
pattern (positions, types and identifiers) through which it factors,
together with the induced affine map.
Concretely, reduce it to its support. Delete zero-slope blocks; retain
an equation when its two slopes are independent. A nonzero rank-one block
has slopes $(v,0),(0,v)$ or $(v,v)$ and retains the first, second or third
variable, respectively. In the third case adjust the intercept, using the third bit
$z_i=b_{e_i}t_0+x_i+y_i$:
\[
 v(x_i+y_i)=v z_i+b_{e_i}v t_0.
\]
For a singleton block, retain it exactly when its slope is nonzero.
The key records positions, question types and identifiers, and the full
reduced affine map, including its intercept. Pulling a projected table
back to $U$ and reducing again gives the same key: the third-bit intercept
correction occurs twice and cancels. Such keys therefore already occur
among tables on equation tuples.

Take vertices to be key orbits under addition of a constant from $K$.
Choose the representative whose intercept has $K$-component zero. Restore
a table's output by adding its removed constant to the vertex label.
Every labeling then defines functions satisfying
\begin{equation}\label{eq:fold-share}
 F_Q(P+h e_0^\top)=F_Q(P)+h,\qquad
 F_U(P'\pi)=F_O(P').
\end{equation}
Here $e_0^\top$ is the first-coordinate functional.

\paragraph{The test and completeness.}
Sample $U\sim\omega^{\otimes k}$, a uniform table $P$, noise $a\sim\mu$, and a uniform
$l\in\F_2^m$, where $m=1+2k$. Require
\begin{equation}\label{eq:matrix-test}
 F_U(P)=F_U(P+a l^\top).
\end{equation}
Equality of restored outputs is a translation constraint on the two vertex
labels. The sampling probabilities are its rational weights. Call this
weighted translation game $\mathcal H$. Loops and repeated edges are allowed
for now.

Given a global assignment satisfying at least $1-\xi$ of the equations,
answer each retained variable by its global bit, and each retained equation
by the global first two bits and the required third bit. Label a key
representative by evaluating its table on this answer and applying $C$.
Equivariance makes restoration consistent. With probability at least
$1-k\xi$, every equation in $U$ is satisfied. All support answers are then
projections of one legal vector $z_U$. The compared values are
\[
 C(Pz_U),\qquad C(Pz_U+a(l^\top z_U)).
\]
Since $z_U\ne0$, the first argument is uniform and $l^\top z_U$ is an
independent fair bit. Consequently
\begin{equation}\label{eq:completeness}
 \val(\mathcal H)\ge1-k\xi-p_*/2.
\end{equation}

\paragraph{Why a gadget is needed.}
Without one, take $\mathcal V=K$, $C=\mathrm{id}$ and $a$ uniform on $K$.
For a legal answer $z$, nonzero since $t_0=1$, the evaluation labeling
$F(P)=Pz$ is folded and passes \eqref{eq:matrix-test} exactly when
$a(l^\top z)=0$. As $l^\top z$ is a fair bit independent of $a$, it
accepts with probability exactly $\frac12+\frac1{2|K|}$. The gadget's
nonlinear $C$ absorbs the noise, so acceptance approaches one by
\eqref{eq:completeness}, while every linear family of large restriction
rank still detects it. Soundness exploits this detection through
Lemma~\ref{lem:noise-comparison}.

\section{The latent noise gadget}
\label{latent:section}

We prove the contract \eqref{eq:gadget-contract}: for every $p_*>0$ we
construct a threshold $r_*$, and then, for every sufficiently large $\ell$
divisible by $3$, finite data
$K\le\mathcal V$, $C$ and a rational law $\mu$. Here $x$ is uniform on
$\mathcal V$, $a$ has
law $\mu$, and the two are independent. All vector spaces and linear maps
in this section are over $\F_2$ unless another field is specified.

\subsection{One quadratic block}

Take $F=\F_q$, where $q=2^d$, and put $X=B_{\rm blk}=F^3$ and $d_{\rm blk}=3d$.
Identify their binary duals with $F^3$ using the nondegenerate pairing
$\Tr(u\cdot v)$, where $\Tr:F\to\F_2$ is the field trace. Define
\[
 Q(x)=(x_2x_3,x_1x_3,x_1x_2),\qquad
 D(x,a)=Q(x+a)+Q(x)+Q(a),\qquad
 C_{\rm blk}(x,y)=y+Q(x).
\]
For each one-dimensional $F$-subspace $A=Fv$ of $F^3$, let
\[
 D_A=\{w:v\cdot w=0\},\qquad
 U_A=\{(a,Q(a)+w):a\in A,\ w\in D_A\}.
\]
Expanding one coordinate, $(x_2+a_2)(x_3+a_3)+x_2x_3+a_2a_3=x_2a_3+x_3a_2$,
shows, by symmetry, that $D(x,a)=x\times a$ is the cross product; signs do not matter in
characteristic two. Hence
\[
 v\cdot D(x,tv)=t\,v\cdot(x\times v)=0,\qquad D(tv,t'v)=tt'\,v\times v=0.
\]
The second identity says that $Q|_A$ is additive, so $U_A$ is a binary subspace of dimension $d_{\rm blk}$.
The trace annihilator of $D_A$ is $A$: containment holds, and both have
binary dimension $d$. The spaces $U_A$ together span $X\oplus B_{\rm blk}$.
Indeed, the coordinate lines supply all horizontal vectors, since $Q$
vanishes there, and their vertical planes span $B_{\rm blk}$.

The cancellation identity is
\[
 C_{\rm blk}((x,y)+(a,Q(a)+w))+C_{\rm blk}(x,y)
 =w+D(x,a)\in D_A.
\]
For a uniform update from $U_A$, this difference is uniform on $D_A$,
even with the input and $a$ fixed. There are $q^2+q+1$ field lines.
Set $\theta=(q^2+q+1)^{-1}$. Removing the zero update, which has
probability $q^{-3}$ and never changes the output, gives
\begin{equation}
 \Pr[\text{output changes under a uniform nonzero update in }U_A]
 =\frac{1-q^{-2}}{1-q^{-3}}=1-\theta.
 \label{latent:block-error}
\end{equation}

Now fix a binary subspace $S\le B_{\rm blk}^*$ and a linear extension
$z\mapsto(\gamma_z,z)\in X^*\oplus B_{\rm blk}^*$, before choosing $A$.
Write $\gamma_z(x)=\Tr(g(z)\cdot x)$, where $g:S\to F^3$ is binary-linear.
For $z\ne0$,
\begin{equation}
 (\gamma_z,z)|_{U_A}=0
 \quad\Longleftrightarrow\quad
 A=Fz\ \text{ and }\ z\cdot g(z)=\sqrt{z_1z_2z_3}.
 \label{latent:kernel}
\end{equation}
The vertical updates first force $z\in A$. The remaining condition is
\[
 \Tr\bigl(t\,g(z)\cdot z+t^2z_1z_2z_3\bigr)=0
 \quad\text{for every }t\in F.
\]
Since $\Tr(t^2c)=\Tr(t\sqrt c)$, nondegeneracy of the trace proves
\eqref{latent:kernel}. Square roots here are unique because squaring
is an automorphism of a finite field of characteristic two.

\subsection{Generic spaces work for every extension}

Call an $r$-dimensional binary subspace $S\le F^3$ \emph{generic} if no
two distinct nonzero elements are $F$-proportional and, for every
binary-linear $g:S\to F^3$, the alignment equation in
\eqref{latent:kernel} has at most $3r$ nonzero solutions.
For each fixed $r$, a uniform $r$-dimensional subspace is generic with
probability tending to one as $q$ tends to infinity. We prove this
uniform statement, including its quantification over all $g$.

Use independent uniform parameters $y_{ij}\in F$, and write
\[
 z_i(t)=\sum_{j=1}^r y_{ij}^2t_j,\qquad
 c_t=\prod_{i=1}^3\left(\sum_{j=1}^r y_{ij}t_j\right),
 \qquad t\in\F_2^r.
\]
Thus $c_t=\sqrt{z_1(t)z_2(t)z_3(t)}$. Squaring preserves uniformity.
The map $t\mapsto z(t)$ fails to be injective with probability at most
$(2^r-1)/q^3$; conditional on injectivity its image is a uniform
$r$-space. Two distinct nonzero binary masks are linearly independent
over $F$, so their images are independent uniform vectors. Their
probability of being nonzero and proportional is $O(q^{-2})$.
A finite union over pairs proves the first genericity condition.

Fix a set $E\subseteq\F_2^r\setminus\{0\}$ with $|E|>3r$, and form
\[
 M_{t,(i,j)}=t_jz_i(t)\qquad(t\in E,\ 1\le i\le3,\ 1\le j\le r).
\]
An extension aligning at all masks in $E$ exists exactly when
$c\in\operatorname{col}(M)$: the unknowns are the three coordinates of
the $r$ values defining $g$. This single system handles every observer.

First regard the parameters as indeterminates, and put
\[
 \mathbb L_0=\F_2(y_{ij}^2)\subset \mathbb L_1=\F_2(y_{ij}).
\]

\begin{lemma}[A cubic outside a quadratic span]\label{lem:cubic}
If $E\subseteq\F_2^r\setminus\{0\}$ and $|E|>3r$, then $c$ is not in the
column space of $M$, even over an
algebraic closure of $\mathbb L_1$.
\end{lemma}

\begin{proof}
The squarefree monomials in all the variables $y_{ij}$ form a basis of
$\mathbb L_1$ over $\mathbb L_0$. In particular, the monomials
$y_{1j}y_{2k}y_{3l}$ are linearly independent over $\mathbb L_0$.
To see this, clear denominators. Coefficients then have only even
exponents, whereas these monomials have distinct exponent-parity
patterns, so cancellation is impossible. The choice of the smaller
coefficient field is essential.

Suppose that $c$ belongs to the column space of $M$ over such a closure.
Every row annihilator of $M$ with entries in
$\mathbb L_0$ then annihilates $c$. Independence of the preceding monomials
forces it to annihilate each vector $(t_jt_kt_l)_{t\in E}$ separately.
By the annihilator characterization of a column space, all these
vectors belong to $\operatorname{col}_{\mathbb L_0}(M)$.
Let $V_j$ be the $\mathbb L_0$-span of evaluations on $E$ of positive-degree
monomials of degree at most $j$. Repeated indices cover degrees one
and two, giving
\[
 V_3\subseteq\operatorname{col}_{\mathbb L_0}(M)\subseteq V_2\subseteq V_3.
\]
Multiplication by any input coordinate takes $V_2$ into $V_3$, so it
preserves this common space. The space therefore contains every
positive-degree monomial. It contains the indicator of each point of
$E$: the usual Boolean point-indicator polynomial has a positive
coordinate factor because the point is nonzero. Consequently the
column space has dimension $|E|>3r$, a contradiction.
\end{proof}

\noindent Now specialize.
Let $\varrho_M$ be the formal rank of $M$. There is a nonzero augmented
$(\varrho_M+1)$-minor of $[M\mid c]$, whereas every unaugmented minor of that
size vanishes identically. Its polynomial degree is at most
$2\varrho_M+3\le6r+3$. Thus any specialization admitting an alignment on $E$
zeros this polynomial, including specializations where the rank drops.
A nonzero polynomial vanishes at independent uniform field elements
with probability at most its total degree divided by $q$: induct on the
number of variables, using the leading coefficient and the
one-variable root bound. Union over sets $E$ of size $3r+1$ now
finishes the proof. Their number depends only on $r$; if no such sets
exist the condition is automatic. Conditioning on injectivity costs a
factor tending to one.

Here is the rank estimate we will use. Put $H_0=0$ and
$H_r=\sum_{i=1}^r1/i$, and let
\[
 r_A=\dim\operatorname{im}\bigl[z\mapsto(\gamma_z,z)|_{U_A}\bigr],
 \qquad r=\dim S.
\]
For generic $S$, \eqref{latent:kernel} gives
\begin{equation}
 r-r_A\le1,\qquad
 \Pr_A[r_A<r]\le3r\theta,\qquad
 \E_A(H_r-H_{r_A})\le3\theta.
 \label{latent:good-loss}
\end{equation}
The first bound uses the absence of proportional pairs. For arbitrary
$S$ we still have
\begin{equation}
 \Pr_A[r_A<r]\le2^r\theta,\qquad H_r-H_{r_A}\le H_r.
 \label{latent:arbitrary-loss}
\end{equation}

\subsection{Recursion reduces the nonlinear error}

Construct spaces $R_t\le P_t$, surjections $\lambda_t:R_t\to B_{\rm blk}$, and
maps $C_t:P_t\to B_{\rm blk}$ satisfying
\begin{equation}
 C_t(u+h)=C_t(u)+\lambda_t(h)\qquad(h\in R_t).
 \label{latent:recursive-symmetry}
\end{equation}
At height zero all spaces are $B_{\rm blk}$ and both maps are the identity.
At each subsequent height use one child for each pair in
\[
 \mathcal I=\{(A,J):J:B_{\rm blk}\xrightarrow{\sim}U_A\},
 \qquad
 P_t=\prod_{\mathcal I}P_{t-1},
 \qquad
 C_t(u)=C_{\rm blk}\left(\sum_{(A,J)}J C_{t-1}(u_{A,J})\right).
\]
All isomorphisms here are binary. On the product of child translation
spaces define the surjections
\[
 L(h)=(\lambda_{t-1}(h_i))_i,\qquad
 \mathcal A((b_i)_i)=\sum_iJ_i b_i
   =(\mathcal A_X((b_i)_i),\mathcal A_B((b_i)_i)).
\]
The second map is onto because the update spaces span $X\oplus B_{\rm blk}$.
Set
\[
 R_t=\ker(\mathcal A_XL),\qquad
 \lambda_t=(\mathcal A_BL)|_{R_t}.
\]
This last map is onto. Such translations change only the vertical
input of the block, proving \eqref{latent:recursive-symmetry}.
In particular, a uniform input has uniform output.

The noise chooses a uniform child at every level and adds an independent
uniform element of $B_{\rm blk}$ at the selected leaf. Each line has the same
number of isomorphisms, so a uniform child can be exposed as a uniform
line followed by a fresh uniform isomorphism.

We need uniformity even after conditioning on a child's output change.
For any fixed child perturbation $a$, write
$\Delta_a(u)=C_{t-1}(u+a)+C_{t-1}(u)$. Then
\[
 (\Delta_a,C_{t-1})(u+h)
 =(\Delta_a(u),C_{t-1}(u)+\lambda_{t-1}(h)).
\]
Translation preserves the uniform input law. Surjectivity of
$\lambda_{t-1}$ therefore makes the child output uniform conditional
on $\Delta_a$ and $a$. The other child outputs remain independent and
uniform. Since $\mathcal A$ is onto, the block input is uniform
conditional on the chosen child and its change.

All children are identical copies, so the law of that change is
independent of the chosen child index. Given a nonzero change and the
chosen line, the fresh uniform $J$ sends the change to a uniform
nonzero element of $U_A$. Equation \eqref{latent:block-error} thus
gives the exact recursion
\begin{equation}
 p_0=1-q^{-3},\qquad p_t=(1-q^{-3})(1-\theta)^t,
 \label{latent:recursive-error}
\end{equation}
where $p_t$ is the probability that $C_t$ changes under its noise.

\subsection{Linear rank survives the adaptive path}
\label{latent:rank}

A lift of $S\le B_{\rm blk}^*$ is a linear map $z\mapsto\psi_z\in P_t^*$ with
$\psi_z|_{R_t}=z\lambda_t$. We show that a lift determines a single
extension at the current block, before a child is chosen.
Restrict $\psi_z$ to $\prod_{\mathcal I}R_{t-1}$. It vanishes on
$\ker L$, hence descends to a functional $\chi_z$ on
$\prod_{\mathcal I}B_{\rm blk}$. The functional $\chi_z-z\mathcal A_B$ vanishes
on $\ker\mathcal A_X$. Since $\mathcal A_X$ is onto, there is a unique
$\gamma_z\in X^*$ such that
\begin{equation}
 \chi_z-z\mathcal A_B=\gamma_z\mathcal A_X,\qquad
 \phi_{A,J}(z)=J^*((\gamma_z,z)|_{U_A}).
 \label{latent:common-extension}
\end{equation}
Uniqueness makes $\gamma_z$ linear in $z$, and common to all children.

For $h\in R_{t-1}$ in child $(A,J)$ this gives
$\psi_z(h)=\phi_{A,J}(z)(\lambda_{t-1}h)$. For a chosen child, let
$S'=\operatorname{im}\phi_{A,J}$. Choose a linear section $\iota:S'\to S$ and restrict $\psi_{\iota(z')}$ to that
child's input. This is a lift of $S'$. Repeat along the noise path,
choosing sections by a fixed rule from the revealed data. If a
continuing functional detects the leaf noise, its ancestor detects
the original noise, because that noise is supported only in the
selected child at every level.

Choose $r_0\ge1$ and $0<\zeta_0<1$ with
\[
 H_{r_0}2^{r_0}\zeta_0\le1.
\]
Choose $q$ large enough that $d_{\rm blk}\ge r_0$ and that, at every rank from
$1$ to $r_0$, the proportion of nongeneric subspaces is at most
$\zeta_0$. Start from any generic $r_0$-dimensional space $S_0\le B_{\rm blk}^*$ together
with a lift of it.

Let $\mathcal F_j$ contain the full path history just before choosing
the next line and isomorphism. The current space, lift, and common
extension are $\mathcal F_j$-measurable. After the line $A_j$ is
revealed, the next rank is already fixed:
\[
 \rho_{j+1}=\dim E_j,\qquad
 E_j=\operatorname{im}[z\mapsto(\gamma_z,z)|_{U_{A_j}}].
\]
Only then expose $J_j$. The space $S_{j+1}=J_j^*E_j$ is uniform among
subspaces of dimension $\rho_{j+1}$. Its lift may depend on $J_j$; this
is harmless because genericity holds for every lift on a generic
space.

The rest is a probability statement, which we isolate.

\begin{lemma}[The race]\label{lem:race}
Let $\theta>0$, and let $r_0=\rho_0\ge\rho_1\ge\dots\ge0$ be integers with
$r_0\ge1$ and $B_0,B_1,\ldots\in\{0,1\}$, all adapted to a filtration
$(\mathcal F_j)$. Suppose that
\[
 \E B_j\le\zeta_0,\qquad
 \E[H_{\rho_j}-H_{\rho_{j+1}}\mid\mathcal F_j]\le3\theta+H_{r_0}2^{r_0}\theta B_j,
 \qquad H_{r_0}2^{r_0}\zeta_0\le1.
\]
Then $s=\lfloor H_{r_0}/(8\theta)\rfloor$ satisfies $\Pr[\rho_s>0]\ge\frac12$.
\end{lemma}

\begin{proof}
Taking expectations, each step lowers $\E H_{\rho_j}$ by at most $4\theta$.
So $\E H_{\rho_s}\ge H_{r_0}-4\theta s\ge H_{r_0}/2$, while
$H_{\rho_s}\le H_{r_0}\ind_{\{\rho_s>0\}}$.
\end{proof}

\noindent The bound on $\E B_j$ is unconditional. The exceptional term is
removed only after averaging over the history, not under conditioning on it.

To apply the lemma, let $B_j$ indicate a positive-rank nongeneric current
space. Then $B_0=0$ and $\E[B_{j+1}\mid\mathcal F_j,A_j]\le\zeta_0$,
because $S_{j+1}$ is uniform of its dimension. The conditional loss bound
follows from \eqref{latent:good-loss}--\eqref{latent:arbitrary-loss}.
So $\rho_s>0$ with probability at least $1/2$.
The uniform leaf symbol has not yet been exposed. A family of rank
$\rho_s>0$ detects it with probability $1-2^{-\rho_s}\ge1/2$.
Ancestor detection therefore shows that every lift of a generic
$r_0$-dimensional space detects the noise with probability at least $1/4$.

\subsection{A random root and one quotient}

Fix a binary basis of $F^3$ and put $K=B_{\rm blk}=\F_2^{\ell}$. The
following lemma completes the gadget; its random root automorphism plays
the role of the fresh isomorphisms $J$ at every level.

\begin{lemma}[Root step]\label{lem:root}
Let $r_0$, $\zeta_0$, $d$ and $s$ be as in Section~\ref{latent:rank}, so
that (i)~every lift on $P_s$ of a generic $r_0$-dimensional subspace of
$B_{\rm blk}^*$ detects the depth-$s$ noise with probability at least
$1/4$, and (ii)~a uniform $r_0$-dimensional subspace of $B_{\rm blk}^*$
is generic with probability at least $1-\zeta_0\ge1/2$. Then, with output
dimension $\ell=d_{\rm blk}=3d$, there are a finite binary space
$\mathcal V$ containing $K$, a map $C:\mathcal V\to K$ and a rational law
$\mu$ on $\mathcal V$ such that:
(a)~$C(x+h)=C(x)+h$ for all $x\in\mathcal V$ and $h\in K$;
(b)~for uniform $x$ and independent $a\sim\mu$,
$\Pr[C(x+a)\ne C(x)]=p_s$, the nonlinear error of the depth-$s$
construction; and
(c)~every $G\le\mathcal V^*$ with $\dim(G|_K)\ge r_0$ satisfies
$\Pr_{a\sim\mu}[\exists g\in G:g(a)\ne0]\ge1/8$.
\end{lemma}

\begin{proof}
Here (i) is the conclusion of Section~\ref{latent:rank}; in (ii),
$\zeta_0\le1/2$ because $H_{r_0}2^{r_0}\zeta_0\le1$ and $r_0\ge1$.
Take one copy of the depth-$s$
construction for every binary automorphism $\tau\in GL(B_{\rm blk})$:
\[
 P=\prod_{\tau}P_s,\qquad R=\prod_{\tau}R_s,\qquad
 \widetilde C((u_\tau))=\sum_{\tau}\tau\,C_s(u_\tau),\qquad
 \Lambda((h_\tau))=\sum_{\tau}\tau\,\lambda_s(h_\tau).
\]
By \eqref{latent:recursive-symmetry},
$\widetilde C(u+h)=\widetilde C(u)+\Lambda(h)$ for $h\in R$, and
$\Lambda:R\to K$ is onto because $\lambda_s$ and each $\tau$ are. Take the
single quotient
\[
 \mathcal V=P/\ker\Lambda,\qquad R/\ker\Lambda\simeq K.
\]
Then $\widetilde C$ descends to $C:\mathcal V\to K$ with $C(x+h)=C(x)+h$
for $h\in K$, and $K$ is an actual subspace of $\mathcal V$. The noise
draws a uniform $\tau$ and, independently of $\tau$, a depth-$s$ noise
sample $\xi_s\in P_s$; it places $\xi_s$ in copy $\tau$ and passes to the quotient. A uniform input
on $P$ pushes forward to a uniform input on $\mathcal V$, since all cosets of
$\ker\Lambda$ have the same size. Since $\tau$ is bijective, the output
change event is preserved exactly, and its probability is still $p_s$.
This proves (a) and (b); $\mu$ is rational, being a pushforward of
uniform choices on finite sets.

For (c), fix $G\le\mathcal V^*$ with $\dim(G|_K)\ge r_0$. It is fixed
before the noise, and so before $\tau$. Put $S=G|_K$ and fix an
$r_0$-dimensional subspace $T\le S$. Let $\varpi:P\to\mathcal V$ be the
quotient map; on $R$ it is $\Lambda$ under the identification
$R/\ker\Lambda\simeq K$. Thus for $g\in G$ and $h\in R_s$ in copy $\tau$,
$g(\varpi\,\iota_\tau h)=(g|_K)\tau\lambda_s(h)$, where $\iota_\tau$
includes factor $\tau$ into $P$. The map $g\mapsto(g|_K)\tau$ sends $G$
onto $\tau^*S$, which contains $\tau^*T=\{z\tau:z\in T\}$. For uniform
$\tau$, the space $\tau^*T$ is uniform among the $r_0$-dimensional
subspaces of $B_{\rm blk}^*$, since $GL(B_{\rm blk})$ acts transitively
on them. By (ii) it is generic with probability at least $1-\zeta_0$.

Fix such a $\tau$ with $\tau^*T$ generic, and a linear section
$\varsigma:\tau^*T\to G$ of $g\mapsto(g|_K)\tau$. The functionals
$\psi_z=\varsigma(z)\circ\varpi\circ\iota_\tau$ satisfy
$\psi_z|_{R_s}=z\lambda_s$, so they form a lift of $\tau^*T$. Since
$a=\varpi\,\iota_\tau\xi_s$ with $\xi_s$ drawn independently of $\tau$, by
(i) this lift detects $\xi_s$ with probability at least $1/4$. Each
$\psi_z$ is the pullback of $\varsigma(z)\in G$, so
$\psi_z(\xi_s)=\varsigma(z)(a)$, and $G$ then detects $a$. Averaging over $\tau$
gives detection probability at least $(1-\zeta_0)/4\ge1/8$.
\end{proof}

\paragraph{Parameter order and finite data.}
Choose $r_0$ large enough that $e^{1-H_{r_0}/8}\le p_*$; this is
possible because the harmonic sums diverge. Set $r_*=r_0$ and choose
$\zeta_0$. The field degree $d$ may then be any integer large enough that
$3d\ge r_0$ and that, at every rank from $1$ to $r_0$, the proportion of
nongeneric subspaces is at most $\zeta_0$. Finally set $\ell=d_{\rm blk}=3d$,
$\theta$ and $s$:
\[
 p_*\ \longrightarrow\ r_0=r_*\ \longrightarrow\ \zeta_0\
 \longrightarrow\ d\ \text{(any large enough)}\ \longrightarrow\ \ell=3d,\ \theta,\ s.
\]
Equation \eqref{latent:recursive-error} gives
\[
 p_s\le e^{-\theta s}\le e^{\theta-H_{r_0}/8}\le p_*.
\]
Lemma~\ref{lem:root} now gives the contract
\eqref{eq:gadget-contract} with $r_*=r_0$, for every sufficiently large
$\ell$ divisible by $3$.


\section{From high acceptance to private decoders}\label{sec:decoders}

Fix $p_*$ and its gadget threshold $r_*$. Set $\eta=2^{-r_*}/16$.
Apply Input 2 at retention $\eta/4$ to obtain $\alpha,r_{\rm s}$. Choose
$\ell$ divisible by $3$ and large enough for that input and the gadget, and so that
\begin{equation}\label{eq:matrix-parameters}
 2^{-(\ell-r_{\rm s})}<\alpha/8,\qquad 2^{-\ell}<\eta/4.
\end{equation}
Now fix the gadget and splitting $\mathcal V=K\oplus W$. Only $k$ grows
from now on. Write a table as $P=(M,T)$ with values in $K\oplus W$.

Extraction concerns the following experiment.

\paragraph{The precise experiment.}
Sample $U\sim\omega^{\otimes k}$. Independently in each position, retain
the equation with probability $1-\beta$, or a uniform one of its variables
with probability $\beta=k^{-2/3}$. This gives $O$ and $\pi:L_U\to L_O$.
Sample independent uniform maps
\[
 A:K\to\F_2^{r_{\rm s}},\qquad T_O:L_O\to W,\qquad M_O:L_O\to K,
 \qquad S_O=AM_O.
\]
The messages are exactly
\begin{equation}\label{eq:messages}
 (U,A,T_O\pi,S_O\pi),\qquad (O,A,T_O,S_O).
\end{equation}
Neither player receives $M_O$ or $\pi$. Each may use the fixed function
family and independent private randomness, but only its own message.

\begin{proposition}[Extraction]\label{prop:extraction}
There are $k_0$ and $\gamma>0$, depending only on the fixed auxiliary data,
such that the following holds for every parity instance and every cube
$k\ge k_0$. If a labeling of $\mathcal H$ passes \eqref{eq:matrix-test} with
probability at least $99/100$, then there are such strategies, returning
legal answers, with
\begin{equation}\label{eq:agreement}
 \Pr[\pi a_U=a_O]\ge\gamma.
\end{equation}
The strategies need not be efficient.
\end{proposition}

\noindent The rest of this section proves the proposition. Choose $k_0$
so that every $k\ge k_0$, with $m=1+2k$, gives $2^{-m}\le\eta/4$, makes
$\ell+m$ exceed Input~2's ambient-dimension threshold at retention
$\eta/4$, and makes $v_k\le\gamma_1$ in \eqref{eq:tv} and
\eqref{eq:gamma1} below; this is possible because $v_k\to0$ as
$k\beta^2=k^{-1/3}\to0$.

\subsection{Replacing the noise and extracting constant outputs}

For a function on a finite binary space, put
$\widehat f(\phi)=\E_x f(x)(-1)^{\phi(x)}$.
A nonzero character averages to zero: translation by a point where it is
one reverses its sign. Thus characters form an orthonormal basis, Fourier
inversion holds, and Parseval gives
$\sum_\phi|\widehat f(\phi)|^2=\E f^2$. Also
\[
 \ind_{u=v}=\E_{\sigma\in K^*}(-1)^{\sigma(u+v)}.
\]
For a noise law $\nu$ on $\mathcal V$ and a function $F$ from tables
$L_U\to\mathcal V$ to $K$, let $\mathrm{acc}_\nu(F)=\Pr[F(P)=F(P+al^\top)]$,
with $P$ and $l$ uniform and $a\sim\nu$. Let $\nu_K$ be uniform on $K$.

\begin{lemma}[Noise comparison]\label{lem:noise-comparison}
$\mathrm{acc}_{\nu_K}(F)\ge2^{-r_*}\bigl(8\,\mathrm{acc}_\mu(F)-7\bigr)$.
\end{lemma}

\begin{proof}
Expand the equality indicator over $\sigma\in K^*$, and each
$(-1)^{\sigma F}$ in characters $g=(g_1,\ldots,g_m)\in(\mathcal V^*)^m$ of
the table space. Orthogonality over $P$ and averaging over $l$ give
\[
 \mathrm{acc}_\nu(F)=\E_\sigma\sum_g|\widehat{F_\sigma}(g)|^2\lambda_\nu(g),
 \qquad
 \lambda_\nu(g)=\Pr_{a\sim\nu}[g_i(a)=0\text{ for all }i]\ge0,
\]
and the squared coefficients sum to one for each $\sigma$. Call $g$ low when
the restrictions $g_i|_K$ span dimension less than $r_*$, and let
$m_{\rm low}$ be the average mass on low frequencies. For the others, the family
$\Span\{g_i\}$ has restriction to $K$ of dimension at least $r_*$, and
some element of it is nonzero at $a$ exactly when some $g_i$ is. So detection gives
$\lambda_\mu\le7/8$. Hence
\[
 \mathrm{acc}_\mu\le m_{\rm low}+\tfrac78(1-m_{\rm low}),\qquad\text{that is,}\qquad
 m_{\rm low}\ge8\,\mathrm{acc}_\mu-7.
\]
On the low frequencies
$\lambda_{\nu_K}=2^{-\dim\Span(g_i|_K)}\ge2^{-r_*}$.
\end{proof}

The bound is linear, so it averages over $U$. Test acceptance $0.99$
therefore gives average uniform-$K$ acceptance at least
$0.92\cdot2^{-r_*}=14.72\eta$. This noise leaves $T$ fixed. Since
acceptance at fixed $(U,T)$ lies in $[0,1]$, a fraction at least
$(14.72\eta-4\eta)/(1-4\eta)>10\eta$ of these pairs satisfies
\begin{equation}\label{eq:good-pair}
 \Pr_{M,a\in K,l}[F_U(M,T)=F_U(M+a l^\top,T)]\ge4\eta.
\end{equation}

\paragraph{A constant-output slice.}
The \emph{step retention} of a set of matrices is the probability that the
uniform rank-one step $M\mapsto M+al^\top$, zero step included, keeps a
uniform member inside it. For any function on matrices, its acceptance is
the fiber-size-weighted average of the step retentions of its fibers.

\begin{lemma}[Constant-output slice]\label{lem:slice}
Let $\ell$, and then $\ell+m$, exceed the thresholds of Input~2 at retention
$\eta/4$, and let $2^{-\ell},2^{-m}\le\eta/4$. Let $f$ be a function on
$\Hom(\F_2^m,K)$ with values in an arbitrary set, and suppose that the fiber
of a value $u$ has step retention at least $\eta$. Then there are a row
map $A_*:K\to\F_2^{r_{\rm s}}$, a value $S_*$, and at most $r_{\rm s}$ column
conditions $Mq_j=t_j$ that cut out a nonempty slice of the row fiber
$A_*M=S_*$ on which $f=u$ with density at least $\alpha$.
\end{lemma}

\begin{proof}
The step is zero with probability
$p_{\rm zero}=2^{-\ell}+2^{-m}-2^{-\ell-m}\le\eta/2$ by hypothesis.
Conditioning on nonzero factors leaves retention at least $\eta/2$.

Identify $M$ with its graph
\[
 \Gamma_M=\{(x,M^\top x):x\in\F_2^\ell\},\qquad
 \dim(\Gamma_M\cap \Gamma_N)=\ell-\rank(M-N).
\]
These form the matrix chart in the Grassmann graph on
$\F_2^\ell\oplus\F_2^m$. A chart vertex has
$(2^\ell-1)(2^m-1)$ chart neighbors and
$(2^\ell-1)(2^{m+1}-2)$ neighbors in total. For the latter count, choose
a hyperplane in the vertex and a different one-dimensional extension in
the quotient. For the former, use the unique nonzero rank-one
factorization over $\F_2$. Exactly half the neighbors are in the chart.
The lifted fiber has Grassmann retention at least $\eta/4$.

Input 2 gives relative density $\alpha$ in an interval. Intersecting it
with the chart cannot lower density, since the entire fiber is already
in the chart. A basis $(d_i,s_i)$ of $A_0$ and a basis $(t_j,q_j)$ of
$B_0^\perp$, for the dot product on $\F_2^\ell\oplus\F_2^m$ under which
$\Gamma_M^\perp=\{(Mv,v)\}$, express the interval's intersection with the
chart as the
nonempty slice
\begin{equation}\label{eq:slice}
 \{M:d_i^\top M=s_i^\top,\ Mq_j=t_j\},
 \qquad i\le n_{\rm row},\ j\le n_c,\ n_{\rm row}+n_c\le r_{\rm s}.
\end{equation}
On at least an $\alpha$ fraction of it, $f$ equals $u$. Its row
conditions say $A_*M=S_*$, where $A_*$ has the rows $d_i^\top$ padded with
zeros, and there are at most $r_{\rm s}$ column conditions.
\end{proof}

The lemma needs no folding assumption, and values outside the fiber play
no role. The constant output, rather than a general affine target, is all
that is used below.

\paragraph{Good row advice occurs often.}
For fixed $(U,T)$ write $f(M)=F_U(M,T)$. Call a nonempty row fiber $AM=S$
good if at most $r_{\rm s}$ column conditions cut out a nonempty slice of it
on which one value of $f$ has density at least $\alpha$.

\begin{lemma}[Good row advice is frequent]\label{lem:advice}
Assume the dimension hypotheses of Lemma~\ref{lem:slice}. If $(U,T)$
satisfies \eqref{eq:good-pair}, the union $\mathcal R$ of all good
row fibers, over all choices of $A$ and $S$, has measure at least $\eta$.
\end{lemma}

\begin{proof}
Suppose that $\mathcal R$ has measure less than $\eta$. Give each matrix in
$\mathcal R$ its own new value outside $K$, and call the result $f'$.
A rank-one step can lose a match only if an endpoint lies in $\mathcal R$.
Both endpoints are uniform, so $f'$ has acceptance at least
$4\eta-2\eta=2\eta$. Each new value has a one-point fiber, whose step
retention is exactly $p_{\rm zero}\le\eta/2$. If every fiber of a value in $K$ had
step retention below $\eta$, the averaging identity would give acceptance
below $\eta+p_{\rm zero}<2\eta$. So some $u\in K$ has a fiber of step retention at
least $\eta$ under $f'$. Lemma~\ref{lem:slice} gives a row fiber
$A_*M=S_*$ and a slice of it, cut out by at most $r_{\rm s}$ column
conditions, on which $f'=u$ with density at least $\alpha$. Wherever
$f'=u$, also $f=u$. So that row fiber is good, and lies in $\mathcal R$.
But on $\mathcal R$ the function $f'$ takes no value in $K$, a
contradiction.
\end{proof}

For each matrix in $\mathcal R$, at least one of the $2^{\ell r_{\rm s}}$ possible row maps
witnesses goodness. Averaging also over the good pairs gives, under a
uniform full table and independent uniform $A$,
\begin{equation}\label{eq:g0}
 \Pr[(U,A,T,AM)\text{ is good row advice}]
 \ge g_0:=10\eta^2 2^{-\ell r_{\rm s}}.
\end{equation}
Select a certificate deterministically from each good advice, using a
fixed order and the function family. Put $Z=\Span\{q_j\}$. If every
vector of $Z$ has first bit zero, translations $M\mapsto M+h e_0^\top$
with $h\in\ker A$ preserve the slice and translate its output by $h$.
Every value then has density at most
$2^{-\dim\ker A}\le2^{-(\ell-r_{\rm s})}<\alpha/8$, a contradiction. Therefore
\begin{equation}\label{eq:answer-list}
 \mathcal A=Z\cap\{v:e_0^\top v=1\},\qquad 1\le|\mathcal A|\le2^{r_{\rm s}}.
\end{equation}
The first player samples uniformly from this list, and gives any legal
answer on advice that is not good. This uses only its own message.

\subsection{Transferring a positive expectation}

Fix $U$. Let $\mathcal L_{\rm pad}$ be the law of the padded full table
$(M_O\pi,T_O\pi)$, with its intercept, after averaging over the projection
choices, and let $\mathcal L_{\rm unif}$ be the law of a uniform table on
$L_U$.

\begin{lemma}[Padded tables look uniform]\label{lem:tv}
With $n=|\mathcal V|$, the total variation distance satisfies
\begin{equation}\label{eq:tv}
 v_k=\bigl\|\mathcal L_{\rm pad}-\mathcal L_{\rm unif}\bigr\|_{\rm TV}
 \le\frac12\sqrt{\Bigl(1+\frac{n-1}{3}\beta^2\Bigr)^k-1}.
\end{equation}
The same bound holds jointly with an independent uniform $A$.
\end{lemma}

\begin{proof}
Blocks are independent. A retained equation has uniform slopes. In a
singleton block, the slope law $Q_0$ is the equally weighted mixture of
$(a,0),(0,a),(a,a)$ for uniform $a\in\mathcal V$. Let $P_0$ be uniform on
$\mathcal V^2$. The ratio $Q_0/P_0$ is $n$ at zero, $n/3$ at the $3(n-1)$
other points of these three sets, and zero elsewhere, so
$\chi^2(Q_0\Vert P_0):=\E_{P_0}(Q_0/P_0-1)^2=(n-1)/3$. Mixing with weight
$\beta$ multiplies the centered ratio by $\beta$, and $1+\chi^2$ multiplies
over independent blocks. Cauchy--Schwarz bounds total variation by half
the square root of $\chi^2$. The intercept is covered too: an independent
uniform intercept stays uniform after adding third-bit corrections that
depend on the slopes.
\end{proof}

\noindent The lemma does \emph{not} compare joint laws that include $\pi$.
Since $k\beta^2=k^{-1/3}$, the bound tends to zero.

Under the unrestricted uniform law, let $J$ mean that row advice is good
and the chosen column conditions hold; let $E$ mean that the output equals
the chosen constant. The conditions impose at most $\ell r_{\rm s}$ extra scalar
equations on a nonempty row fiber. Hence
\[
 \Pr(J)\ge g_0 2^{-\ell r_{\rm s}},\qquad
 \E[\ind_J(\ind_E-\alpha/2)]\ge(\alpha g_0/2)2^{-\ell r_{\rm s}}.
\]
This score depends only on the full table, $U,A$, and the deterministic
certificate selected from its row advice. Its range has length one, so
\eqref{eq:tv} changes its expectation by at most $v_k\le\gamma_1$. Hence,
in the projected experiment its expectation is at least
\begin{equation}\label{eq:gamma1}
 \gamma_1:=(\alpha g_0/4)2^{-\ell r_{\rm s}}>0.
\end{equation}

Only now condition on $D=(U,O,\pi,A,T_O,S_O)$. The first player's
certificate is fixed, while $M_O$ is uniform on its row fiber. By sharing,
the selected output equality is an equality for $F_O$. Let $p_D$ be the
probability of its projected column slice and $\rho_D$ the constant-output
density there; take them to be zero when no such slice exists. The
conditional score is $p_D(\rho_D-\alpha/2)$. It is nonpositive unless the
slice is nonempty and $\rho_D\ge\alpha/2$, and is always at most one.
Consequently these useful data occur with probability at least $\gamma_1$.
No witness has been sent to the second player.

\subsection{A decoder that does not know the certificate}

The second player answers with a general decoding rule. Let $H\le K$,
$d_H=\dim H$ and $X=\Hom(L_O,H)$, and let $g:X\to K$ be folded:
$g(N+he_0^\top)=g(N)+h$ for $h\in H$. Identify
\[
 X^*=H^*\otimes L_O,\quad
 \langle\varphi\otimes v,N\rangle=\varphi(Nv),\quad
 (\varphi\otimes v)(h')=\varphi(h')v.
\]
The decoder chooses $\sigma\in K^*$ uniformly, puts
$f_\sigma=(-1)^{\sigma g}$, and samples $\Phi\in X^*$ with probability
$|\widehat f_\sigma(\Phi)|^2$; Parseval normalizes this law. If
$\sigma|_H\ne0$, it returns $\Phi(h_\sigma)$ for a fixed $h_\sigma\in H$ with
$\sigma(h_\sigma)=1$. Otherwise it returns a fixed legal answer.
In plain terms, a dense constant-output slice forces Fourier mass onto
frequencies supported on the column directions that define the slice,
and folding turns such a frequency into a legal answer.

\begin{lemma}[Private decoding]\label{lem:decode}
The decoder always returns a legal answer. Let $Z'\le L_O$ have dimension
$s$. If on some nonempty slice $N|_{Z'}=\kappa$ one value of $g$ has density
$\rho\ge2^{-d_H}$, then with probability at least
$2^{-d_Hs}(\rho-2^{-d_H})^2$ the decoder returns an answer in $Z'$.
\end{lemma}

\begin{proof}
Folding gives $f_\sigma(N+h_\sigma e_0^\top)=-f_\sigma(N)$, so every
sampled frequency satisfies $\langle\Phi,h_\sigma e_0^\top\rangle=1$, that
is, $\Phi(h_\sigma)$ has first coordinate one. Let
$C_\sigma=\E[f_\sigma\mid N|_{Z'}=\kappa]$. The character identity for the
dense value gives $\E_\sigma|C_\sigma|\ge\rho$. The characters vanishing on
$H$ form a $2^{-d_H}$ fraction, so
$\E_\sigma[\ind_{\sigma|_H\ne0}|C_\sigma|]\ge\rho-2^{-d_H}$. The maps
vanishing on $Z'$ have annihilator $\mathcal F_{\rm col}=H^*\otimes Z'$, of
size $2^{d_Hs}$. Averaging a character over the affine slice retains
precisely these frequencies, so for any point $N_*$ of the slice
\[
 C_\sigma=\sum_{\Phi\in\mathcal F_{\rm col}}\widehat f_\sigma(\Phi)
                (-1)^{\langle\Phi,N_*\rangle}.
\]
Two applications of Cauchy--Schwarz give
\[
 \Pr[\sigma|_H\ne0,\ \Phi\in\mathcal F_{\rm col}]
 =\E_\sigma\Bigl[\ind_{\sigma|_H\ne0}
   \sum_{\Phi\in\mathcal F_{\rm col}}|\widehat f_\sigma(\Phi)|^2\Bigr]
 \ge2^{-d_Hs}\bigl(\rho-2^{-d_H}\bigr)^2.
\]
On this event $\Phi$ maps $H$ into $Z'$, so the answer lies in $Z'$.
\end{proof}

The second player fixes a representative $M_0$ of its row fiber, with
$AM_0=S_O$, and applies the decoder to $g(N)=F_O(M_0+N,T_O)$ with $H=\ker A$, using only its
own message; $g$ is folded by \eqref{eq:fold-share}. For useful data $D$
put $Z'=\pi Z$, so $s\le r_{\rm s}$. Since the slice is nonempty,
$\kappa(\pi q_j):=t_j-M_0\pi q_j$ defines a linear map $Z'\to H$. The
transferred slice $N|_{Z'}=\kappa$ has a constant output with density
$\rho=\rho_D\ge\alpha/2$; this description is
used only in the analysis. Since
$2^{-d_H}\le2^{-(\ell-r_{\rm s})}<\alpha/8$ and $d_Hs\le\ell r_{\rm s}$,
Lemma~\ref{lem:decode} gives probability at least
$2^{-\ell r_{\rm s}}(\alpha/4)^2$ of an answer in $Z'$.
On this event, $a_O\in Z'$ and is legal. Because $\pi$ preserves the first
coordinate, $a_O\in\pi\mathcal A$. The first player's independent uniform
choice from \eqref{eq:answer-list} hits a suitable preimage with probability
at least $2^{-r_{\rm s}}$. Averaging over useful data proves \eqref{eq:agreement}
with the fixed positive constant
\begin{equation}\label{eq:gamma}
 \gamma=\gamma_1 2^{-\ell r_{\rm s}}(\alpha/4)^2 2^{-r_{\rm s}}.
\end{equation}
Both strategies use only their own messages. The agreement bound concerns
the original unconditioned experiment.


\section{Clean coordinates force disagreement}\label{sec:clean}

\begin{proposition}[Clean coordinates]\label{prop:clean}
If $\OPT(\mathcal E)\le4/5$, in particular in a NO instance with $\xi<1/100$,
\emph{every} pair of strategies that use only their own
messages \eqref{eq:messages} and private randomness, and return legal
answers, satisfies
\[
 \Pr[\pi a_U=a_O]\le
 \exp\Bigl[-\frac{2^{-\dim W-r_{\rm s}}}{3600}k^{1/3}\Bigr].
\]
\end{proposition}

\begin{proof}
By \eqref{eq:base-gap}, the equation--variable game has value at most
$14/15$.

Fix $A$ and combine the visible advice as
\[
 D_O=(T_O,S_O),\qquad E_A=W\oplus\operatorname{im}A.
\]
Conditional on questions and projections, its original intercept and all
coefficient columns are independent uniform elements of $E_A$. Their law
does not depend on equation identifiers or selected variable positions.
Call a position clean when it is a singleton and its advice coefficient
is zero. Conditional on $A$, positions are independently clean with
probability
\[
 p_A=\beta 2^{-\dim W-\rank A},\qquad
 \Pr[e_i=e,j_i=j,i\text{ clean}\mid A]
       =\omega_e\frac{\beta}{3|E_A|}.
\]
Conditioning on cleanliness thus preserves exactly the base game's
question-pair law.

Let $I$ be the clean set. Fix only the following background: $A$, $I$,
the \emph{original intercept on $O$}, and all questions, projection choices,
and coefficients outside $I$. Do not fix a certificate, output, success
event, or answer randomness. For each background of positive probability,
the clean question pairs remain independent samples of the base game.
This follows from the displayed identity, independence across positions,
and independence of the original intercept. Repeated identifiers do not
change the product law; positions have distinct local answer coordinates.

Each player can reconstruct its own full message from the background
and its own clean questions. Indeed, pulling a singleton coefficient
$d_i$ back to an equation gives slopes
$(d_i,0),(0,d_i)$ or $(d_i,d_i)$. The padded intercept is
\[
 c_U=c_O+\sum_{\substack{i\text{ singleton}\\j_i=3}}b_{e_i}d_i.
\]
At a clean position $d_i=0$, so both slopes and the intercept correction
vanish. The equation player can reconstruct without knowing which variable
was selected; the projected player can reconstruct without knowing the
clean equations. Each runs its original strategy and keeps its clean
answers. These are legal strategies for the $|I|$-fold base game. Fixing
background merely selects a pair of strategies for this bound; it adds
no message to the original experiment. Private randomness is a mixture
of deterministic strategies and cannot increase the optimum.

Full agreement implies success in every clean position. Input~3 and
\eqref{eq:base-gap} give
\[
 \Pr[\pi a_U=a_O\mid\text{background}]\le(1-1/3600)^{|I|}.
\]
For $I=\varnothing$ use the bound one. We do not assume independence of
the strategies' successes. Since $|I|$ is binomial conditional on $A$,
\begin{align}
 \Pr[\pi a_U=a_O]
 &\le\E_A(1-p_A/3600)^k\nonumber\\
 &\le\left(1-\frac{\beta 2^{-\dim W-r_{\rm s}}}{3600}\right)^k
 \le\exp\left[-\frac{2^{-\dim W-r_{\rm s}}}{3600}k^{1/3}\right]
 \longrightarrow0.\label{eq:clean-upper}
\end{align}
\end{proof}
This bounds the \emph{same unconditioned experiment} as
\eqref{eq:agreement}. Once the auxiliary spaces are fixed, sufficiently
large $k$ gives a contradiction. Thus, when $\xi<1/100$, the intermediate
game has the gap
\begin{equation}\label{eq:intermediate-gap}
 \mathrm{YES}:\val(\mathcal H)\ge1-k\xi-p_*/2,
 \qquad \mathrm{NO}:\val(\mathcal H)\le99/100.
\end{equation}

\section{Removing weights, amplifying, and writing the output}\label{sec:finish}

\paragraph{Rounding weights.}
If there are $N_c$ weighted constraint occurrences, let $Q=\lceil N_c/\tau\rceil$.
Start with $\lfloor Qw_e\rfloor$ copies of occurrence $e$ and distribute
the remaining copies to the largest fractional parts. Then
\[
 \sum_e n_e=Q,\qquad |w_e-n_e/Q|\le1/Q,\qquad
 \tfrac12\sum_e|w_e-n_e/Q|\le\tau.
\]
The weight differences sum to zero. Hence the weight of every subset of
constraints changes by at most $\tau$, uniformly over labelings. The
optimum changes by at most $\tau$ as well. No large common denominator
is needed.

\paragraph{A simple bipartite graph.}
Replace every copy by a path of four edges with three fresh internal
vertices. The first three permutations are identities; the fourth is
the original constraint. For fixed endpoint labels, all four edges can
be satisfied exactly when the original constraint holds. Otherwise exactly
three can be satisfied. Distinct paths have disjoint interiors, so their
optima can be achieved simultaneously. The resulting game $G_{\rm bip}$ satisfies
\begin{equation}\label{eq:subdivision}
 \left|\val(G_{\rm bip})-\left(1-\frac{1-\val(\mathcal H)}4\right)\right|
 \le\tau/4.
\end{equation}
Place original vertices and path midpoints on one side and the other
internal vertices on the other. A loop becomes a four-cycle, and parallel
copies have distinct interiors. The graph is simple and bipartite.
Reversing an edge leaves a binary translation unchanged.

\paragraph{Parameter order.}
Choose positive rationals $\eps_0\le\eps$ and $\delta_0\le\delta$.
If $\tau\le1/200$, a NO instance has gap at least $1/800$ after subdivision.
Choose a fixed integer $t$ such that
\begin{equation}\label{eq:final-t}
 \left(1-\frac1{10\,240\,000}\right)^t\le\delta_0.
\end{equation}
This choice precedes the alphabet: the bound in \eqref{eq:repetition}
does not depend on it. The denominator is $16\cdot800^2$.
Next choose rational budgets
\[
 p_*\le\eps_0/t,\qquad \tau\le\min\{1/200,\eps_0/t\}.
\]
The remaining order is
\[
 p_*\ \longrightarrow\ r_*\ \longrightarrow\ \eta
 \ \longrightarrow\ \alpha,r_{\rm s},\ell_0
 \ \longrightarrow\ \ell=3d\ge\ell_0,\ \mathcal V,K,W
 \ \longrightarrow\ k,\beta\ \longrightarrow\ \xi.
\]
Here $\ell_0$ is the largest of Input~2's threshold for $\ell$ at retention
$\eta/4$, the threshold for \eqref{eq:matrix-parameters}, and $3d_0$, where
$d_0$ is such that every $d\ge d_0$ meets the gadget's genericity and
$3d\ge r_0$ requirements; it exists because the proportion of nongeneric
subspaces tends to zero.
Choose a cube $k\ge k_0$, with $k_0$ from Proposition~\ref{prop:extraction},
so large that the bound \eqref{eq:clean-upper} is below $\gamma$; this is
possible because $k\beta=k^{1/3}\to\infty$. Then the extraction and clean
coordinate bounds contradict each other in the NO case. As in the source,
$k$ is a cube so that $\beta=k^{-2/3}$ is rational; this matters only for an
effective search for the constants, which the theorem does not need.
All thresholds are uniform over the instance and function family.
Finally choose a positive rational
\[
 \xi<\min\{1/100,\eps_0/(2kt)\}.
\]
For fixed target errors, these finite constants can be hardcoded into one
machine for every input length; their existence therefore suffices here.
We make no uniform computability claim in the error targets.

\paragraph{The final game.}
In the YES case put $\mathrm{err}=k\xi+p_*/2\le\eps_0/t$.
Output $G=G_{\rm bip}^{\otimes t}$. A good labeling in every coordinate and a union
bound yield
\[
 \val(G)\ge1-t(\mathrm{err}+\tau)/4\ge1-\eps_0/2\ge1-\eps.
\]
In the NO case, \eqref{eq:repetition} and \eqref{eq:final-t} give
\[
 \val(G)\le\left(1-\frac1{10\,240\,000}\right)^t\le\delta_0\le\delta.
\]
The product graph is simple: a pair of product vertices specifies at most
one edge in each coordinate. It is bipartite, and independent uniform
base edges give uniform product edges. Its constraints are translations
over the fixed alphabet $K^t=\F_2^{\ell t}$.

\paragraph{Polynomial output length.}
All dimensions and repetition counts are fixed before the input. If the
processed parity system has $N_{\mathcal E}$ occurrences, there are
$N_{\mathcal E}^k$ tuples. Each has a fixed number of tables and noise
outcomes. Each constraint weight is a product of a fixed number of
rationals of polynomial bit length. Keys are computed by fixed-dimensional
linear algebra; even comparing every pair of keys is polynomial.
The rounding multiplicity is linear in the number of tests, and the final
power is fixed. The permutation tables have fixed size relative to the
input. Thus the entire explicit game can be written in polynomial time,
completing the argument for Theorem~\ref{thm:main}.

\paragraph{Attribution and changes of presentation.}
The construction and main proof are from \cite{UGT}. Its matrix-chart
argument credits Barak--Kothari--Steurer \cite{BKS}; the needed
calculation is proved here. The table matches passages of this note with
those of \cite{UGT}; passages not listed are inherited in condensed form.

\begin{center}
\footnotesize
\begin{tabular}{@{}>{\raggedright\arraybackslash}p{0.2\textwidth}
 >{\raggedright\arraybackslash}p{0.18\textwidth}
 >{\raggedright\arraybackslash}p{0.56\textwidth}@{}}
\toprule
This note & Source & Relation\\
\midrule
\S3.1 polar map; \S3.2 genericity &
 \S3.1, Lemma 3.2; \S3.2, Lemma 3.6 &
 Same statements. The polar map is written as a cross product; a square
 parametrization replaces square roots of monomials.\\
\S3.5, root step (Lemma~\ref{lem:root}) & \S3.5 &
 Replaced, proved here: copies indexed by automorphisms, one quotient,
 threshold $r_0$, output dimensions $3d$. Source: copies indexed by
 injections, two quotients, threshold $3d+2$, all large dimensions.\\
Erasure (Lemma~\ref{lem:advice}) & Lemma 5.5 &
 Erased matrices get distinct new values outside $K$, instead of a random
 recoloring in $K$ with concentration and counting bounds.\\
Certificates (Lemmas~\ref{lem:slice}, \ref{lem:decode}; \S\S4.1--4.3) &
 Theorem 5.1, Lemmas 5.5, 5.7, 5.8 &
 Constant output, the case $z=0$ that the proof of Theorem 5.1 already
 gives, instead of a general affine target $Mz+u$.\\
\S5; \S6 & \S6; \S7 & Inherited; \S6 only shows that the constants exist,
 while the source also shows that they are computable.\\
\bottomrule
\end{tabular}
\end{center}

The changes simplify the presentation and some proof steps; none is a
claim of a new theorem or an independent verification of the preprint.

\appendix
\section{Lean check}\label{app:lean}

OpenAI's repository \cite{UGT} contains a Lean~4 formalization of the
source theorem, with a Comparator challenge that permits only the
standard axioms. We report one run of this check, made on 9 October 2026
at the revision pinned in \cite{UGT}, as an unprivileged user on a
four-core x86-64 cloud machine. The check script
(\texttt{ugc\_lean\_check.py}) and its complete output logs
(\texttt{ugc-lean-ugt-check.log} and the directory \texttt{ugc\_out})
are deposited with this note at
\href{https://doi.org/\depositdoi}{\texttt{doi:\depositdoi}}.

\begin{center}
\footnotesize
\begin{tabular}{@{}ll@{}}
Repository & \texttt{openai/math}, revision
 \texttt{fd4aeeb2ee4fc729c18d98444fed42fd0529eeeb}\\
Tools & Lean 4.34.1 (pinned toolchain); Comparator commit \texttt{d03acab}
 (tag \texttt{v4.34.0})\\
Command & \texttt{lake env comparator
 ComparatorChallenges/UniqueGamesTheorem.json}\\
Declaration & \texttt{OAI.UniqueGamesTheorem.theorem11}\\
Axioms & \texttt{propext}, \texttt{Classical.choice},
 \texttt{Quot.sound}\\
Verdict & ``Lean default kernel accepts the solution'', ``Your solution
 is okay!''
\end{tabular}
\end{center}

\paragraph{Procedure and result.} Following the repository's
instructions, our script ran \texttt{lake update} (manifest and toolchain
unchanged), \texttt{lake exe cache get} and the command above in the
directory \texttt{lean}, then a separate \texttt{lake build} of the
solution and challenge modules and \texttt{\#print axioms}. The challenge
module states the declaration's type: \texttt{BinaryGapReduction}~$\eps$~$\delta$
is nonempty for all real $\eps,\delta\in(0,1/2)$. Comparator exited with
code~$0$, the build succeeded, and the solution's axioms include no
\texttt{sorryAx}; only the challenge module uses \texttt{sorry}, by design,
to state the type. The run took 66 minutes.

\paragraph{Caveats.} The sandbox was not enforced: \texttt{landrun}
started, but in our test a write outside its allowed directory succeeded,
and the \texttt{systemd-run} guard recommended in Comparator's README was
unavailable. The check therefore lacked the recommended isolation. We did
not audit whether the formal statement matches Theorem~\ref{thm:main}.
Above all, the run concerns OpenAI's formal development of the source
theorem. It does not check this note's modified root construction
(Lemma~\ref{lem:root}), its other changes listed under Attribution, or
anything else written here.

\footnotesize
\begin{thebibliography}{9}
\setlength{\itemsep}{1pt}
\bibitem{UGT}
OpenAI. \emph{The Unique Games Theorem}. Preprint, 23 September 2026,
58 pages. Pinned source revision:
\href{https://github.com/openai/math/tree/fd4aeeb2ee4fc729c18d98444fed42fd0529eeeb/preprints/The-Unique-Games-Theorem-September-23-2026}{GitHub source and paper}.

\bibitem{Hastad}
Johan H\aa stad. \emph{Some Optimal Inapproximability Results}.
Journal of the ACM 48(4), 798--859, 2001.
Input 1 uses Theorem 5.4 and its proof in the
\href{https://people.kth.se/~johanh/optimalinap.pdf}{author manuscript}.

\bibitem{KMS}
Subhash Khot, Dor Minzer, and Muli Safra.
\emph{Pseudorandom Sets in Grassmann Graph Have Near-Perfect Expansion}.
Annals of Mathematics 198(1), 1--92, 2023.
Input 2 uses Theorem 1.12 in the
\href{https://eccc.weizmann.ac.il/report/2018/006/revision/2/download}{second 2018 revision}.

\bibitem{DS}
Irit Dinur and David Steurer. \emph{Analytical Approach to Parallel
Repetition}. Proceedings of STOC, 624--633, 2014.
Input 3 uses Corollary 1.2 in the
\href{https://www.dsteurer.org/paper/productgames.pdf}{10 June 2014 author version}.

\bibitem{BKS}
Boaz Barak, Pravesh K. Kothari, and David Steurer.
\emph{Small-Set Expansion in Shortcode Graph and the 2-to-2 Conjecture}.
Proceedings of ITCS, LIPIcs 124, 9:1--9:12, 2019.

\bibitem{Fortnow}
Lance Fortnow (prepared with Claude). \emph{The Unique Games Theorem: An
Expository Account of the Proof}. 7 October 2026,
\href{https://github.com/fortnow/AI-Papers}{github.com/fortnow/AI-Papers}.
\raggedright
\bibitem{ShortProof}
Ariel Elboim. \emph{The Unique Games Theorem: A Short Self-Contained
Proof}. Preprint, 2026,
\href{https://doi.org/10.5281/zenodo.23265572}{doi:10.5281/zenodo.23265572}.
\end{thebibliography}
\end{document}
